\documentclass[11pt,a4paper]{article}

\usepackage[margin=1in]{geometry}
\usepackage{amsmath}
\usepackage{booktabs}
\usepackage{graphicx}
\usepackage[hidelinks]{hyperref}

\title{Explainable Pneumonia Classification from Chest X-Rays Using DenseNet-121, CBAM, and Multimodal Fusion}
\author{Kshitiz Reddy.M \\ M.Tech Artificial Intelligence \\ Bennett University}
\date{\today}

\begin{document}
\maketitle

\begin{abstract}
This work evaluates DenseNet-121 for binary pneumonia classification from
pediatric chest X-rays, with Convolutional Block Attention Module (CBAM)
attention and post-hoc Grad-CAM explanations. On a 624-image held-out test
split, the vision model achieved an AUROC of 0.9604 (95\% CI: 0.9412--0.9757)
and an AUPR of 0.9628 (95\% CI: 0.9377--0.9825). We additionally evaluate a
late-fusion configuration using synthetic tabular attributes. Its improved
metrics demonstrate that the implementation can exploit correlated multimodal
signal, not clinical validity with real patient records. We report all
thresholded and ranking metrics with bootstrap confidence intervals and use
lung segmentation masks to quantify the fraction of Grad-CAM activation inside
the lung field.
\end{abstract}

\section{Introduction}
Pneumonia remains an important use case for computer-aided interpretation of
chest radiographs. Deep neural networks can achieve high classification
performance, but visual explanations must be interpreted cautiously: a
plausible heatmap does not establish clinically valid reasoning. This work
therefore combines classification evaluation with Grad-CAM-based localization
measurement.

\section{Related Work}
CBAM applies sequential channel and spatial attention with limited additional
overhead~\cite{woo2018cbam}. Recent pneumonia studies have combined CBAM,
DenseNet-121, and Grad-CAM, although their task definitions and data splits
differ from this binary study~\cite{dey2026cbam,potharaju2025enhanced,
rashid2025framework}. Grad-CAM is widely used for weakly supervised
localization~\cite{selvaraju2017gradcam,shahi2025weakly}. DenseNet variants
have also been evaluated directly on the Kermany chest X-ray data set
\cite{bundea2024pneumonia,erukude2025explainable}.

\section{Methods}
\subsection{Architecture}
An ImageNet-pretrained DenseNet-121 backbone produces a 1024-dimensional image
representation. CBAM is inserted after feature extraction and before the
classification head. For the fusion configuration, four tabular inputs are
encoded into a 64-dimensional representation and concatenated with the image
representation before classification.

\subsection{Explainability}
Grad-CAM~\cite{selvaraju2017gradcam} produces class-specific heatmaps from the
final convolutional layer. A pretrained lung-segmentation model provides lung
masks. The lung-energy fraction is the fraction of heatmap activation contained
inside the segmented lung region. This is a quantitative proxy for anatomical
localization, not a clinical validation of the explanation.

\section{Dataset and Preprocessing}
We use the pediatric Chest X-Ray Pneumonia data set~\cite{kermany2018identifying},
containing 5,856 images (1,583 normal and 4,273 pneumonia). The original
624-image test split is retained. The training and validation pools are
re-split using the best available patient identifiers; this reduces leakage
risk but cannot prove complete patient-level separation for normal images,
which lack equivalent identifiers. Images are resized to $224\times224$,
replicated from grayscale to RGB, and normalized. Training images receive
horizontal-flip and small-rotation augmentation.

% Paste these revised sections into the main .tex file supplied by Claude.
% Required preamble packages: \usepackage{booktabs,graphicx,amsmath}

\section{Results}
\label{sec:results}

Table~\ref{tab:final-metrics} reports performance on the held-out test split
($n=624$) at a pre-specified probability threshold of 0.5. Confidence intervals
are percentile-bootstrap 95\% intervals from 2,000 resamples of the test set.
The vision model uses DenseNet-121 with CBAM. The fusion model uses the same
image branch and additionally receives synthetic tabular features; consequently,
the fusion result demonstrates the implementation's ability to use correlated
tabular signal and must not be interpreted as clinical validation.

\begin{table}[htbp]
\centering
\caption{Held-out test performance. Values in brackets are bootstrap 95\%
confidence intervals.}
\label{tab:final-metrics}
\small
\begin{tabular}{lcc}
\toprule
Metric & Vision + CBAM & Fusion + CBAM \\
\midrule
Accuracy  & 0.8638 [0.8365, 0.8895] & 0.8702 [0.8429, 0.8959] \\
Precision & 0.8224 [0.7876, 0.8566] & 0.8280 [0.7927, 0.8627] \\
Recall    & 0.9974 [0.9920, 1.0000] & 1.0000 [1.0000, 1.0000] \\
F1-score  & 0.9015 [0.8800, 0.9217] & 0.9059 [0.8843, 0.9263] \\
AUROC     & 0.9604 [0.9412, 0.9757] & 0.9899 [0.9810, 0.9961] \\
AUPR      & 0.9628 [0.9377, 0.9825] & 0.9921 [0.9839, 0.9978] \\
\bottomrule
\end{tabular}
\end{table}

The fusion configuration has higher point estimates for every reported metric.
Its AUROC increases by 0.0295 and its AUPR by 0.0293 relative to the vision
configuration. The AUROC and AUPR intervals do not overlap; this descriptive
observation is not a formal paired hypothesis test. Threshold-based intervals
for accuracy and F1-score overlap, so the results do not support a strong claim
of improvement for those metrics at the fixed threshold.

\section{Corrections to Experimental Setup}

\subsection{Evaluation Metrics and Uncertainty}
Models are evaluated using accuracy, precision, recall, F1-score, AUROC, and
AUPR. Accuracy, precision, recall, and F1-score are calculated at a fixed
threshold of 0.5; this threshold is not tuned on the held-out test split. AUROC
and AUPR summarize ranking performance across thresholds. For every metric, a
percentile-bootstrap 95\% confidence interval is calculated using 2,000
resamples of the held-out test set with replacement.

\subsection{Data Splitting}
The original 624-image test split is preserved. The original training and
validation pools are merged and then split into training and validation sets
using the best available patient identifiers. Pneumonia filenames encode a
patient identifier; normal-class images do not provide an equivalent identifier
and are therefore treated as unique patients. This reduces, but cannot prove
the complete absence of, patient-level leakage. The resulting experimental
splits contain 4,434 training, 798 validation, and 624 test images.

\subsection{Statistical Replication}
Selected architecture and attention-consistency comparisons are repeated over
three random seeds (42, 123, and 2024) and summarized as mean $\pm$ standard
deviation. Any associated $t$-tests are exploratory because $n=3$ offers low
statistical power; their p-values should not be treated as definitive evidence.

% Replace the existing Proposed row in Table 1 with this row. Do not use the
% earlier 0.9608 value: the evaluated checkpoint has AUROC 0.9604.
% \textbf{Proposed} & \textbf{DenseNet121 + CBAM} & \textbf{Binary CXR}
% & \textbf{86.38\%} & \textbf{82.24\%} & \textbf{99.74\%}
% & \textbf{90.15\%} & \textbf{0.9628} & \textbf{0.9604} \\


\section{Limitations}
The data set is pediatric and from a single public source; performance may not
generalize across hospitals, imaging devices, age groups, or disease
prevalences. The fusion metadata are synthetic and deliberately label-correlated,
so they must not be interpreted as real clinical evidence. Grad-CAM and the
lung-energy fraction are explanatory proxies rather than ground-truth lesion
localization measures. Three-seed comparisons are exploratory because of their
small sample size.

\section{Conclusion}
DenseNet-121 with CBAM provides strong held-out ranking performance for this
research prototype. The fusion configuration improves the reported point
estimates, but its synthetic metadata prevent clinical conclusions. Future work
should validate on external, patient-linked data and compare explanations with
expert annotations.

\bibliographystyle{plain}
\bibliography{references}
\end{document}
